\section{モデルの定式化}

\subsection{半マルコフ過程}

半マルコフ過程は，遷移先は現在状態のみに依存するが，各状態での
滞在時間には指数分布に限らない任意の分布を許す確率過程である．
本稿では，潜伏時間に Pareto 分布を直接導入するため，この枠組みを
用いる．

\subsection{攻撃フェーズと防疫フェーズ}

\begin{definition}[攻撃フェーズ]
状態空間を $\{A_0,A_1,A_2,A_3\}$ とする．
$A_0$ は無感染，$A_1$ は潜伏感染，$A_2$ は発症中，
$A_3$ は破壊完了を表す．攻撃到着は強度 $\lambda>0$ の
Poisson 過程であり，攻撃到着時に $A_0\to A_1$ が起こる．
潜伏時間 $T_d$ の経過後に $A_1\to A_2$，さらに
$T_e\sim\Exp(\mu_e)$ の後に $A_2\to A_3$ が起こる．
\end{definition}

\begin{assumption}[Pareto 潜伏時間]
$T_d\sim\Pareto(t_m,\alpha(I))$ とし，
\[
  \Prob(T_d>t)=\left(\frac{t_m}{t}\right)^{\alpha(I)},\qquad
  f_{T_d}(t)=\alpha(I)t_m^{\alpha(I)}t^{-(\alpha(I)+1)}
\]
for $t\ge t_m>0$ とする．尾指数 $\alpha(I)>0$ は投資 $I$ の関数である．
\end{assumption}

\begin{definition}[防疫フェーズ]
状態空間を $\{D_0,D_1,D_2,D_3\}$ とする．
$D_0$ は通常，$D_1$ は検知，$D_2$ は隔離，$D_3$ は復旧中を表す．
早期検知時刻は $T_{\mathrm{det}}\sim\Exp(\beta(I))$ で，
$T_{\mathrm{det}}<T_d$ の場合に早期検知が生起する．発症後には
確率 $\eta\in(0,1]$ で症状検知される．
\end{definition}

\begin{assumption}[投資応答関数]
$\alpha(I)$ と $\beta(I)$ は $[0,\infty)$ 上の $C^2$ 関数で，
単調増加かつ凹型である．また
\[
  \lim_{I\to\infty}\alpha(I)=\alpha_{\max}<\infty,\qquad
  \lim_{I\to\infty}\beta(I)=\beta_{\max}<\infty,
\]
かつ $\alpha(0)>0,\ \beta(0)>0$ とする．
\end{assumption}

合成状態空間は
\[
S=\{S_N,S_L,S_E,S_S,S_I,S_R,S_F\}
\]
で表す．$S_N$ は正常，$S_L$ は潜伏・未検知，$S_E$ は早期検知済，
$S_S$ は発症後・症状検知待ち，$S_I$ は隔離済，$S_R$ は復旧中，
$S_F$ は全損を表す．

\begin{figure}[htbp]
\centering
\begin{tikzpicture}[
  node distance=16mm and 22mm,
  state/.style={draw,rounded corners,minimum width=13mm,minimum height=8mm,align=center},
  loss/.style={draw,rounded corners,fill=gray!20,minimum width=13mm,minimum height=8mm,align=center},
  arr/.style={-{Latex[length=2mm]},thick}
]
\node[state] (SN) {$S_N$};
\node[state,right=of SN] (SL) {$S_L$};
\node[state,above right=of SL] (SE) {$S_E$};
\node[state,right=of SE] (SI) {$S_I$};
\node[state,right=of SI] (SR) {$S_R$};
\node[state,below right=of SL] (SS) {$S_S$};
\node[loss,right=of SS] (SF) {$S_F$};
\draw[arr] (SN) -- node[above] {$\lambda$} (SL);
\draw[arr] (SL) -- node[left] {$\beta(I)$} (SE);
\draw[arr] (SE) -- node[above] {隔離} (SI);
\draw[arr] (SI) -- node[above] {復旧} (SR);
\draw[arr] (SL) -- node[left] {$T_d$ 後} (SS);
\draw[arr] (SS) -- node[above] {$1-\eta$} (SF);
\draw[arr] (SS) -- node[left] {$\eta$} (SI);
\draw[arr] (SR) to[bend left=35] node[above] {完了} (SN);
\end{tikzpicture}
\caption{合成半マルコフ状態遷移図．}
\end{figure}

\subsection{バックアップと復旧}

\begin{definition}[バックアップスケジュール]
間隔 $\Delta>0$ で $n\in\mathbb{Z}_+$ 世代のバックアップを保持する．
侵入後に取得されたバックアップは汚染されるとし，汚染世代数は
早期検知時に $k=\lceil T_{\mathrm{det}}/\Delta\rceil$，
症状検知時に $k=\lceil T_d/\Delta\rceil$ とする．
$k\ge n$ なら全損，$k<n$ なら復旧可能である．
\end{definition}

\begin{definition}[コンテナベース復旧時間]
復旧は隔離，クローン，再デプロイの三段階から成る．
$T_{\mathrm{iso}}\sim\Exp(\mu_{\mathrm{iso}})$，
$T_{\mathrm{clone}}\sim\Exp(\mu_c)$，
$T_{\mathrm{deploy}}\sim\Exp(\mu_d)$ とし，
\[
  \tau_R:=\mu_{\mathrm{iso}}^{-1}+\mu_c^{-1}+\mu_d^{-1}.
\]
全停止時間は $T_{\mathrm{down}}=T_{\mathrm{iso}}+T_{\mathrm{rec}}$，
$T_{\mathrm{rec}}=T_{\mathrm{scan}}+T_{\mathrm{clone}}+T_{\mathrm{deploy}}$，
$T_{\mathrm{scan}}=\kappa_sT_{\mathrm{det}}$ とする．
\end{definition}

\subsection{期待損失と費用レート}

更新サイクルあたりの期待損失を
\begin{equation}
  \Qfull(\Delta;\alpha)
  =(1-p_e)\left\{\eta(d_1\tau_R+d_0)(1-F)
  +\eta d_1\kappa_s ET+F L\right\}
  \label{eq:qfull}
\end{equation}
と定義する．ここで
\[
  p_e=\Prob(T_{\mathrm{det}}<T_d),\quad
  F(\Delta;\alpha)=\left(\frac{t_m}{n\Delta}\right)^\alpha,\quad
  ET(\Delta;\alpha)=\E[T_d\mathbf{1}_{\{T_d\le n\Delta\}}].
\]
バックアップ固定費，投資費，追加保持費を含む定常費用レートは
\begin{equation}
  C(\Delta,I,n_{\mathrm{ext}})
  =\lambda\Qfull(\Delta;\alpha(I))
  +\frac{c_b}{\Delta}+c_I I+c_n n_{\mathrm{ext}}.
  \label{eq:cost-rate}
\end{equation}

